% TOP_PROBLEM: {"id":30007529,"problem_number":"LOCAL-30007529","title":"Responding to persistent supply shocks","category_id":21,"category":{"id":21,"name":"economics","display_name":"Economics","description":"Open economic theory statements, published research agendas, and proposed scoped specifications; question provenance and historical priority are qualified per record.","slug":"economics","order_index":21,"created_at":"2026-10-08T00:00:00Z"},"status":"provenance_pending","external_url":null,"legacy_ids":[],"published":false,"statement":"Characterize a monetary rule minimizing worst-case stabilization losses when supply-shock persistence and expectation anchoring belong to a stated uncertainty set.","economics":{"original_id":18,"field":"Inflation and monetary policy","kind":"design","formalization_basis":"proposed_specification","source_status":"not_verified","priority_status":"not_established","priority_note":"No explicit primary open-problem statement matching this catalogue item was verified; supporting research is not evidence of first formulation.","earliest_verified_reference":null,"current_reference":{"authors":["Hess Chung","Callum Jones","Antoine Lepetit","Fernando M. Martin"],"title":"Implications of Inflation Dynamics for Monetary Policy Strategies","year":2025,"url":"https://www.federalreserve.gov/econres/feds/files/2025072pap.pdf","doi":"10.17016/FEDS.2025.072","locator":"Abstract, p. 1; Section 1, pp. 2–3","evidence_summary":"The paper studies robust monetary strategies and supply-shock tradeoffs and gives model-based guidance. It does not state the precise minimax persistence-uncertainty problem formulated below as an unresolved conjecture."},"formalization_note":"New robust-control specification informed by current policy research. No exact primary open statement matching it was verified."},"statusQualification":"Provenance pending: an explicit published statement of this exact open question was not verified. This is a catalogue-authored candidate specification, not a certified open conjecture. New robust-control specification informed by current policy research. No exact primary open statement matching it was verified.","statusReviewedAt":"2026-10-08"}
% FIRST_POSTED: 2026-10-08
% ENTRY_KIND: statement_only
% !TeX program = lualatex
\documentclass[11pt]{article}
\usepackage{fontspec}
\usepackage[a4paper,margin=25mm]{geometry}
\usepackage{amsmath,amssymb,mathtools}
\usepackage{xurl}
\usepackage[hidelinks]{hyperref}
\newcommand{\sref}[2]{\hyperlink{source:#1:#2}{[S#2]}}
\setlength{\emergencystretch}{3em}
\begin{document}
% problemId:30007529
\section*{Responding to persistent supply shocks}\label{30007529}
\subsection{Definitions and mathematical statement}
Consider the New Keynesian system
\[
x_t=E_tx_{t+1}-\sigma^{-1}(i_t-E_t\pi_{t+1}-r_t^n),\quad
\pi_t=\beta E_t\pi_{t+1}+\kappa x_t+u_t,\quad
u_{t+1}=\rho u_t+\epsilon_{t+1},
\]
where \(x_t\) is the output gap, \(\pi_t\) inflation relative to target, \(i_t\) the nominal rate, \(r_t^n\) the flexible-price rate, \(\sigma,\kappa>0\), and \(\beta\in(0,1)\). Shock persistence \(\rho\) belongs to a known interval, while the innovation law and possible expectation-anchoring mechanisms belong to a declared compact family \(\mathcal P\). Expectations follow the specified mechanism in each candidate model. The central bank observes noisy indicators and chooses an adapted rule \(r\) satisfying a stated rate bound. Restrict admissible rules to equilibria with finite expected discounted losses.
Characterize the robust rule minimizing
\[
\inf_r\sup_{P\in\mathcal P}E_P\sum_{t\ge0}\beta^t
[\pi_t^2+\lambda x_t^2+\nu(i_t-i_{t-1})^2],
\]
with \(\lambda,\nu\ge0\) fixed policy weights. Determine when learning about persistence justifies accommodation versus stronger stabilization, and how conclusions change when beliefs can become unanchored. Demonstrate computational accuracy and compare robust-policy losses with correctly specified and misspecified benchmarks. This is a model-specific decision problem under explicitly bounded uncertainty; existing qualitative supply-shock guidance does not establish that this precise optimization is an open canonical theorem.

\par\textbf{Formulation and scope.} New robust-control specification informed by current policy research. No exact primary open statement matching it was verified.

\par\textbf{Open-question provenance.} An explicit published open-question statement was not verified. Sources below are supporting literature and must not be treated as a first listing.

\par\textbf{Historical priority.} No explicit primary open-problem statement matching this catalogue item was verified; supporting research is not evidence of first formulation.

\subsection{Short English statement}
Characterize a monetary rule minimizing worst-case stabilization losses when supply-shock persistence and expectation anchoring belong to a stated uncertainty set.

\subsection{Sources}
\begin{itemize}\raggedright
\item[\textnormal{[S1]}]\hypertarget{source:30007529:1}{} Hess Chung, Callum Jones, Antoine Lepetit, Fernando M. Martin. 2025. Implications of Inflation Dynamics for Monetary Policy Strategies. Abstract, p. 1; Section 1, pp. 2–3.
\url{https://www.federalreserve.gov/econres/feds/files/2025072pap.pdf}.
DOI: 10.17016/FEDS.2025.072.
{\small\textit{Supporting or current literature; see evidence qualification. The paper studies robust monetary strategies and supply-shock tradeoffs and gives model-based guidance. It does not state the precise minimax persistence-uncertainty problem formulated below as an unresolved conjecture.}}
\end{itemize}
\end{document}
